\begin{definition}[Pointwise additive inverse of complex sequences]
\label{def:cplx-pointwise-negation}
For a complex sequence $s:\Unary\to\ComplexUp$, the pointwise additive inverse $-s$ is the sequence
$$
(-s)(n):=-(s(n)),
$$
where the endpoint inverse is formed componentwise on the real and imaginary rational-history components of the complex carrier from \autoref{def:complex-carrier}.
\end{definition}
\leandef{BEDC.Derived.ComplexLimitUp.ComplexPointwiseNegation}

\begin{lemma}[Complex distance is invariant under simultaneous negation]
\label{lem:cplx-distance-simultaneous-negation}
For complex histories $a,b$ and a distance history $D$,
$$
\mathsf{CplxDist}(a,b,D)\Longrightarrow \mathsf{CplxDist}(-a,-b,D).
$$
\end{lemma}
\begin{proof}
By \autoref{def:cplx-distance}, the hypothesis is a complex-modulus witness for the componentwise difference $a-b$. The componentwise additive-inverse ledger gives
$$
(-a)-(-b)=-(a-b)
$$
on both rational-history components. In \autoref{def:cplx-modulus}, replacing both components of a difference by their additive inverses does not change the displayed square terms, since rational-history squaring is invariant under additive inverse. The same modulus witness $D$ therefore certifies the distance from $-a$ to $-b$.
\end{proof}
\leanchecked{BEDC.Derived.ComplexLimitUp.ComplexDistance\_simultaneous\_negation}

\begin{theorem}[Complex limit respects pointwise negation]
\label{thm:cplx-lim-pointwise-negation}
If $\mathsf{CplxLim}(s,N,z,M)$ holds, then the pointwise additive-inverse sequence has complex limit $-z$ with the same regularity modulus and the same limit modulus:
$$
\mathsf{CplxLim}(-s,N,-z,M).
$$
\end{theorem}
\begin{proof}
Unpack the limit hypothesis using \autoref{def:cplx-lim}. Its regularity field is a witness for $\mathsf{CplxRegSeq}(s,N)$, its target field carries the complex history $z$, and its modulus field gives the distance estimates from $s(n)$ to $z$ above $M$.

First verify regularity for $-s$ with the unchanged modulus $N$. For a unary precision $k$ and accepted unary indices $n,m$ in the sense of \autoref{def:cplx-regseq}, regularity of $s$ gives a distance witness $D$ from $s(n)$ to $s(m)$ with the required dyadic bound. Applying \autoref{lem:cplx-distance-simultaneous-negation} gives a distance witness with the same $D$ from $(-s)(n)$ to $(-s)(m)$; the numerical bound is unchanged because the distance history has not changed. Thus $\mathsf{CplxRegSeq}(-s,N)$ holds.

The target carrier for $-z$ is supplied by the componentwise additive-inverse closure of the $\ComplexUp$ carrier from \autoref{def:complex-carrier}. Now fix a unary precision $k$ and an accepted unary index $n$ above $M(k)$ as in \autoref{def:cplx-lim}. The original limit witness supplies a distance history $E$ from $s(n)$ to $z$ satisfying the $2^{-k}$ bound. Applying \autoref{lem:cplx-distance-simultaneous-negation} gives the same distance history $E$ from $(-s)(n)$ to $-z$, and the same dyadic estimate remains valid. Since $k$ and $n$ were arbitrary among the accepted indices, $M$ is a limit modulus for $-s$ at $-z$.
\end{proof}
\leanchecked{BEDC.Derived.ComplexLimitUp.ComplexLimit\_pointwise\_negation\_closed}
\leanchecked{BEDC.Derived.ComplexLimitUp.ComplexLimit\_pointwise\_negation}
